\documentclass[11pt,a4paper]{article}
\usepackage[utf8]{inputenc}
\usepackage[T1]{fontenc}
\usepackage[catalan]{babel}
\usepackage{lmodern}
\usepackage[sfdefault,lining]{FiraSans}
\usepackage[a4paper,top=2.3cm,bottom=2.2cm,left=1.8cm,right=1.8cm,headheight=1.3cm,headsep=0.4cm,footskip=1cm]{geometry}
\usepackage{amsmath,amssymb}
\allowdisplaybreaks
\usepackage{icomma}
\usepackage{graphicx}
\usepackage[table]{xcolor}
\usepackage{tikz}
\usepackage{fontawesome5}
\usepackage{ccicons}
\usepackage{fancyhdr}
\usepackage{lastpage}
\usepackage{tasks}
\usepackage{needspace}
\usepackage[most]{tcolorbox}
\usepackage{array}
\usepackage[hidelinks,pdftitle={Fitxa d'exercicis de la unitat 6: Càlcul de derivades. Derivabilitat d'una funció. 2n de Batxillerat MII},pdfsubject={Matemàtiques II},pdfkeywords={derivada, recta tangent, recta normal, regla de la cadena, derivació logarítmica, derivació implícita, derivabilitat, funcions a trossos, paràmetres, PAU}]{hyperref}

\definecolor{navy}{HTML}{1B2748}
\definecolor{accent}{HTML}{E8603A}
\definecolor{teal}{HTML}{1F9E90}
\definecolor{soft}{HTML}{EEF1F8}
\definecolor{muted}{HTML}{5D6685}

\setlength{\parindent}{0pt}
\setlength{\parskip}{3pt}
\settasks{label=\alph*), label-format=\bfseries\color{navy}, label-width=1.3em, item-indent=1.8em, column-sep=0.6em, after-item-skip=5pt, before-skip=3pt, after-skip=4pt}

% ---------- capçalera i peu ----------
\newcommand{\llicencia}{\href{https://creativecommons.org/licenses/by-nc-sa/4.0/deed.ca}{\ccbyncsa}\ \ Llicència Creative Commons Reconeixement-NoComercial-CompartirIgual 4.0 Internacional (CC BY-NC-SA 4.0)}
\pagestyle{fancy}
\fancyhf{}
\fancyhead[L]{\raisebox{-0.35\height}{\includegraphics[height=0.95cm]{logo-gva.png}}}
\fancyhead[C]{\small\color{muted}Fitxa d'exercicis · Unitat 6: Càlcul de derivades · 2n Batxillerat MII}
\fancyhead[R]{\raisebox{-0.35\height}{\includegraphics[height=0.95cm]{logo-vermellar.png}}}
\fancyfoot[L]{\scriptsize\color{muted}\llicencia}
\fancyfoot[R]{\small\color{muted}\thepage}
\renewcommand{\headrulewidth}{0.4pt}
\renewcommand{\headrule}{\color{muted!40}\hrule height \headrulewidth}
\fancypagestyle{primera}{\fancyhf{}\renewcommand{\headrulewidth}{0pt}
  \fancyfoot[L]{\scriptsize\color{muted}\llicencia}
  \fancyfoot[R]{\small\color{muted}\thepage}}

% ---------- nivells i exercicis ----------
\newcommand{\estrelles}[1]{{\color{accent}\ifcase#1\or\faStar\faStar[regular]\faStar[regular]\or\faStar\faStar\faStar[regular]\or\faStar\faStar\faStar\fi}}
\newcounter{exer}
\newcommand{\ex}[1]{\par\needspace{4\baselineskip}\medskip\refstepcounter{exer}\noindent{\color{navy}\bfseries\theexer.}\ \ }

\newcommand{\seccio}[2]{\par\needspace{10\baselineskip}\bigskip
  \begin{tcolorbox}[colback=navy,colframe=navy,coltext=white,arc=2mm,boxrule=0pt,left=3mm,right=3mm,top=1.2mm,bottom=1.2mm]
  {\large\bfseries #1\quad #2}
  \end{tcolorbox}\vspace{-2pt}}


\setlength{\emergencystretch}{3em}

% quadre de recordatori (només a la fitxa d'exercicis)
\newcommand{\recorda}[1]{\par\smallskip\begin{tcolorbox}[colback=accent!7,colframe=accent!55,arc=2mm,boxrule=0.5pt,left=2.5mm,right=2.5mm,top=1mm,bottom=1mm]\small{\color{accent}\faLightbulb[regular]}\ \textbf{Recorda:}\ #1\end{tcolorbox}}
% etiqueta dels exercicis de la PAU d'altres comunitats
\newcommand{\pau}[1]{{\color{teal}\small\bfseries\faUniversity\ PAU · #1.}\ }
\newcommand{\res}[1]{#1}\newcommand{\rest}[1]{#1}

% =====================================================================
% Capçalera FentMates de les fitxes: un color per a cada curs.
% Els colors són els de la pàgina inicial (LEVELS a assets/js/materials.js);
% el logo ve de _eines/logo/pdf/ (python3 _eines/logo/genera.py).
%
% Ús, al final del preàmbul de la unitat:
%   % =====================================================================
% Capçalera FentMates de les fitxes: un color per a cada curs.
% Els colors són els de la pàgina inicial (LEVELS a assets/js/materials.js);
% el logo ve de _eines/logo/pdf/ (python3 _eines/logo/genera.py).
%
% Ús, al final del preàmbul de la unitat:
%   % =====================================================================
% Capçalera FentMates de les fitxes: un color per a cada curs.
% Els colors són els de la pàgina inicial (LEVELS a assets/js/materials.js);
% el logo ve de _eines/logo/pdf/ (python3 _eines/logo/genera.py).
%
% Ús, al final del preàmbul de la unitat:
%   \input{../comu/fentmates.tex}
%   \fmcurs{1eso}
% i, al començament del document:
%   \capcalera{Matemàtiques}{1r ESO}{Fitxa d'exercicis · Unitat 2}{Divisibilitat}
% Opcional abans de \input: \def\fmneed{14} (espai mínim abans de cada secció).
% =====================================================================
\makeatletter
\def\fm@logodir{../../logo/pdf}
\providecommand{\fmneed}{10}
\def\fm@cursnom{}

% clau del curs, color i color del text damunt del color
\newcommand{\fm@curs}[3]{\@namedef{fm@color@#1}{#2}\@namedef{fm@text@#1}{#3}}
\fm@curs{1eso}{E6B84A}{navy}
\fm@curs{2eso}{3B82F6}{white}
\fm@curs{3eso}{FF7A50}{white}
\fm@curs{4eso-a}{A78BFA}{white}
\fm@curs{4eso-b}{6D28D9}{white}
\fm@curs{1batx-mi}{9B1C1C}{white}
\fm@curs{1batx-ccss}{EF6B6B}{white}
\fm@curs{2batx-mii}{1E6B47}{white}
\fm@curs{2batx-ccss}{2BB3A3}{white}

\newcommand{\fmcurs}[1]{%
  \@ifundefined{fm@color@#1}{\PackageError{fentmates}{Curs desconegut: #1}{}}{}%
  \edef\fm@tmp{\noexpand\definecolor{curs}{HTML}{\@nameuse{fm@color@#1}}%
               \noexpand\colorlet{curstext}{\@nameuse{fm@text@#1}}}\fm@tmp
  \colorlet{cursfosc}{curs!72!black}%
  \def\fm@logo{\fm@logodir/vertical-#1.pdf}%
  \def\fm@marca{\fm@logodir/marca-#1.pdf}}

% ---------- capçalera de la primera pàgina ----------
% #1 matèria (petit), #2 curs, #3 tipus de document i unitat, #4 títol
\newcommand{\capcalera}[4]{%
  \gdef\fm@cursnom{#2}%
  \par\noindent
  \begin{tcolorbox}[colback=curs,colframe=curs,boxrule=0pt,arc=2.2mm,width=3.5cm,height=2.15cm,
      valign=center,left=3mm,right=2mm,top=1mm,bottom=1mm,nobeforeafter,box align=center]
    \color{curstext}{\footnotesize #1}\par\vspace{1pt}{\fontsize{19}{21}\selectfont\bfseries #2\par}%
  \end{tcolorbox}\hspace{5mm}%
  \begin{minipage}[c]{\dimexpr\linewidth-3.5cm-5mm-2.9cm\relax}
    {\color{cursfosc}\small\bfseries\MakeUppercase{#3}}\par\vspace{2pt}
    {\color{navy}\fontsize{20}{23}\selectfont\bfseries #4\par}
  \end{minipage}\hfill
  \begin{minipage}[c]{2.6cm}\raggedleft\includegraphics[width=2.6cm]{\fm@logo}\end{minipage}\par
  \vspace{2pt}}

% ---------- barres de secció del color del curs ----------
\renewcommand{\seccio}[2]{\par\needspace{\fmneed\baselineskip}\bigskip
  \begin{tcolorbox}[colback=curs,colframe=curs,coltext=curstext,arc=2mm,boxrule=0pt,left=3mm,right=3mm,top=1.2mm,bottom=1.2mm]
  {\large\bfseries #1\quad #2}
  \end{tcolorbox}\vspace{-2pt}}

% ---------- estrelles de dificultat del color del curs (a Batxillerat no se'n posen) ----------
\@ifundefined{estrelles}{}{%
  \renewcommand{\estrelles}[1]{{\color{curs}\ifcase#1\or\faStar\faStar[regular]\faStar[regular]\or\faStar\faStar\faStar[regular]\or\faStar\faStar\faStar\fi}}}

% ---------- capçalera i peu de les pàgines ----------
% primera pàgina: logos institucionals al peu
\fancypagestyle{primera}{\fancyhf{}\renewcommand{\headrulewidth}{0pt}
  \fancyfoot[L]{\raisebox{-0.32\height}{\includegraphics[height=0.78cm]{logo-gva.png}}\hspace{5mm}%
    \raisebox{-0.32\height}{\includegraphics[height=0.78cm]{logo-vermellar.png}}}
  \fancyfoot[R]{\scriptsize\color{muted}\href{https://creativecommons.org/licenses/by-nc-sa/4.0/deed.ca}{\ccbyncsa}\ \ CC BY-NC-SA 4.0 · fentmates.org\qquad\small\thepage}}
% la resta: el símbol a l'esquerra i el curs a la dreta, amb el títol al centre com abans
\fancyhead[L]{\raisebox{-0.3\height}{\includegraphics[height=0.62cm]{\fm@marca}}}
\fancyhead[R]{\small\bfseries\color{cursfosc}\fm@cursnom}
\renewcommand{\headrule}{\color{curs!45}\hrule height \headrulewidth}
\makeatother

%   \fmcurs{1eso}
% i, al començament del document:
%   \capcalera{Matemàtiques}{1r ESO}{Fitxa d'exercicis · Unitat 2}{Divisibilitat}
% Opcional abans de \input: \def\fmneed{14} (espai mínim abans de cada secció).
% =====================================================================
\makeatletter
\def\fm@logodir{../../logo/pdf}
\providecommand{\fmneed}{10}
\def\fm@cursnom{}

% clau del curs, color i color del text damunt del color
\newcommand{\fm@curs}[3]{\@namedef{fm@color@#1}{#2}\@namedef{fm@text@#1}{#3}}
\fm@curs{1eso}{E6B84A}{navy}
\fm@curs{2eso}{3B82F6}{white}
\fm@curs{3eso}{FF7A50}{white}
\fm@curs{4eso-a}{A78BFA}{white}
\fm@curs{4eso-b}{6D28D9}{white}
\fm@curs{1batx-mi}{9B1C1C}{white}
\fm@curs{1batx-ccss}{EF6B6B}{white}
\fm@curs{2batx-mii}{1E6B47}{white}
\fm@curs{2batx-ccss}{2BB3A3}{white}

\newcommand{\fmcurs}[1]{%
  \@ifundefined{fm@color@#1}{\PackageError{fentmates}{Curs desconegut: #1}{}}{}%
  \edef\fm@tmp{\noexpand\definecolor{curs}{HTML}{\@nameuse{fm@color@#1}}%
               \noexpand\colorlet{curstext}{\@nameuse{fm@text@#1}}}\fm@tmp
  \colorlet{cursfosc}{curs!72!black}%
  \def\fm@logo{\fm@logodir/vertical-#1.pdf}%
  \def\fm@marca{\fm@logodir/marca-#1.pdf}}

% ---------- capçalera de la primera pàgina ----------
% #1 matèria (petit), #2 curs, #3 tipus de document i unitat, #4 títol
\newcommand{\capcalera}[4]{%
  \gdef\fm@cursnom{#2}%
  \par\noindent
  \begin{tcolorbox}[colback=curs,colframe=curs,boxrule=0pt,arc=2.2mm,width=3.5cm,height=2.15cm,
      valign=center,left=3mm,right=2mm,top=1mm,bottom=1mm,nobeforeafter,box align=center]
    \color{curstext}{\footnotesize #1}\par\vspace{1pt}{\fontsize{19}{21}\selectfont\bfseries #2\par}%
  \end{tcolorbox}\hspace{5mm}%
  \begin{minipage}[c]{\dimexpr\linewidth-3.5cm-5mm-2.9cm\relax}
    {\color{cursfosc}\small\bfseries\MakeUppercase{#3}}\par\vspace{2pt}
    {\color{navy}\fontsize{20}{23}\selectfont\bfseries #4\par}
  \end{minipage}\hfill
  \begin{minipage}[c]{2.6cm}\raggedleft\includegraphics[width=2.6cm]{\fm@logo}\end{minipage}\par
  \vspace{2pt}}

% ---------- barres de secció del color del curs ----------
\renewcommand{\seccio}[2]{\par\needspace{\fmneed\baselineskip}\bigskip
  \begin{tcolorbox}[colback=curs,colframe=curs,coltext=curstext,arc=2mm,boxrule=0pt,left=3mm,right=3mm,top=1.2mm,bottom=1.2mm]
  {\large\bfseries #1\quad #2}
  \end{tcolorbox}\vspace{-2pt}}

% ---------- estrelles de dificultat del color del curs (a Batxillerat no se'n posen) ----------
\@ifundefined{estrelles}{}{%
  \renewcommand{\estrelles}[1]{{\color{curs}\ifcase#1\or\faStar\faStar[regular]\faStar[regular]\or\faStar\faStar\faStar[regular]\or\faStar\faStar\faStar\fi}}}

% ---------- capçalera i peu de les pàgines ----------
% primera pàgina: logos institucionals al peu
\fancypagestyle{primera}{\fancyhf{}\renewcommand{\headrulewidth}{0pt}
  \fancyfoot[L]{\raisebox{-0.32\height}{\includegraphics[height=0.78cm]{logo-gva.png}}\hspace{5mm}%
    \raisebox{-0.32\height}{\includegraphics[height=0.78cm]{logo-vermellar.png}}}
  \fancyfoot[R]{\scriptsize\color{muted}\href{https://creativecommons.org/licenses/by-nc-sa/4.0/deed.ca}{\ccbyncsa}\ \ CC BY-NC-SA 4.0 · fentmates.org\qquad\small\thepage}}
% la resta: el símbol a l'esquerra i el curs a la dreta, amb el títol al centre com abans
\fancyhead[L]{\raisebox{-0.3\height}{\includegraphics[height=0.62cm]{\fm@marca}}}
\fancyhead[R]{\small\bfseries\color{cursfosc}\fm@cursnom}
\renewcommand{\headrule}{\color{curs!45}\hrule height \headrulewidth}
\makeatother

%   \fmcurs{1eso}
% i, al començament del document:
%   \capcalera{Matemàtiques}{1r ESO}{Fitxa d'exercicis · Unitat 2}{Divisibilitat}
% Opcional abans de \input: \def\fmneed{14} (espai mínim abans de cada secció).
% =====================================================================
\makeatletter
\def\fm@logodir{../../logo/pdf}
\providecommand{\fmneed}{10}
\def\fm@cursnom{}

% clau del curs, color i color del text damunt del color
\newcommand{\fm@curs}[3]{\@namedef{fm@color@#1}{#2}\@namedef{fm@text@#1}{#3}}
\fm@curs{1eso}{E6B84A}{navy}
\fm@curs{2eso}{3B82F6}{white}
\fm@curs{3eso}{FF7A50}{white}
\fm@curs{4eso-a}{A78BFA}{white}
\fm@curs{4eso-b}{6D28D9}{white}
\fm@curs{1batx-mi}{9B1C1C}{white}
\fm@curs{1batx-ccss}{EF6B6B}{white}
\fm@curs{2batx-mii}{1E6B47}{white}
\fm@curs{2batx-ccss}{2BB3A3}{white}

\newcommand{\fmcurs}[1]{%
  \@ifundefined{fm@color@#1}{\PackageError{fentmates}{Curs desconegut: #1}{}}{}%
  \edef\fm@tmp{\noexpand\definecolor{curs}{HTML}{\@nameuse{fm@color@#1}}%
               \noexpand\colorlet{curstext}{\@nameuse{fm@text@#1}}}\fm@tmp
  \colorlet{cursfosc}{curs!72!black}%
  \def\fm@logo{\fm@logodir/vertical-#1.pdf}%
  \def\fm@marca{\fm@logodir/marca-#1.pdf}}

% ---------- capçalera de la primera pàgina ----------
% #1 matèria (petit), #2 curs, #3 tipus de document i unitat, #4 títol
\newcommand{\capcalera}[4]{%
  \gdef\fm@cursnom{#2}%
  \par\noindent
  \begin{tcolorbox}[colback=curs,colframe=curs,boxrule=0pt,arc=2.2mm,width=3.5cm,height=2.15cm,
      valign=center,left=3mm,right=2mm,top=1mm,bottom=1mm,nobeforeafter,box align=center]
    \color{curstext}{\footnotesize #1}\par\vspace{1pt}{\fontsize{19}{21}\selectfont\bfseries #2\par}%
  \end{tcolorbox}\hspace{5mm}%
  \begin{minipage}[c]{\dimexpr\linewidth-3.5cm-5mm-2.9cm\relax}
    {\color{cursfosc}\small\bfseries\MakeUppercase{#3}}\par\vspace{2pt}
    {\color{navy}\fontsize{20}{23}\selectfont\bfseries #4\par}
  \end{minipage}\hfill
  \begin{minipage}[c]{2.6cm}\raggedleft\includegraphics[width=2.6cm]{\fm@logo}\end{minipage}\par
  \vspace{2pt}}

% ---------- barres de secció del color del curs ----------
\renewcommand{\seccio}[2]{\par\needspace{\fmneed\baselineskip}\bigskip
  \begin{tcolorbox}[colback=curs,colframe=curs,coltext=curstext,arc=2mm,boxrule=0pt,left=3mm,right=3mm,top=1.2mm,bottom=1.2mm]
  {\large\bfseries #1\quad #2}
  \end{tcolorbox}\vspace{-2pt}}

% ---------- estrelles de dificultat del color del curs (a Batxillerat no se'n posen) ----------
\@ifundefined{estrelles}{}{%
  \renewcommand{\estrelles}[1]{{\color{curs}\ifcase#1\or\faStar\faStar[regular]\faStar[regular]\or\faStar\faStar\faStar[regular]\or\faStar\faStar\faStar\fi}}}

% ---------- capçalera i peu de les pàgines ----------
% primera pàgina: logos institucionals al peu
\fancypagestyle{primera}{\fancyhf{}\renewcommand{\headrulewidth}{0pt}
  \fancyfoot[L]{\raisebox{-0.32\height}{\includegraphics[height=0.78cm]{logo-gva.png}}\hspace{5mm}%
    \raisebox{-0.32\height}{\includegraphics[height=0.78cm]{logo-vermellar.png}}}
  \fancyfoot[R]{\scriptsize\color{muted}\href{https://creativecommons.org/licenses/by-nc-sa/4.0/deed.ca}{\ccbyncsa}\ \ CC BY-NC-SA 4.0 · fentmates.org\qquad\small\thepage}}
% la resta: el símbol a l'esquerra i el curs a la dreta, amb el títol al centre com abans
\fancyhead[L]{\raisebox{-0.3\height}{\includegraphics[height=0.62cm]{\fm@marca}}}
\fancyhead[R]{\small\bfseries\color{cursfosc}\fm@cursnom}
\renewcommand{\headrule}{\color{curs!45}\hrule height \headrulewidth}
\makeatother

\fmcurs{2batx-mii}

\begin{document}
\thispagestyle{primera}

\capcalera{Matemàtiques II}{2n Batx}{Fitxa d'exercicis · Unitat 6}{\raggedright Càlcul de derivades. Derivabilitat d'una funció}
\vspace{10pt}
Nom i cognoms: \rule{8.4cm}{0.4pt}\hfill Grup: \rule{1.6cm}{0.4pt}\hfill Data: \rule{2.2cm}{0.4pt}
\vspace{8pt}

\begin{tcolorbox}[colback=soft,colframe=soft,arc=2mm,left=3mm,right=3mm,top=1.5mm,bottom=1.5mm]
\small Resol els exercicis al quadern justificant cada pas. Els exercicis marcats amb {\color{teal}\bfseries\faUniversity\ PAU} són de les proves d'accés a la universitat d'altres comunitats autònomes (de Matemàtiques II i de Matemàtiques aplicades a les CCSS II), traduïts al valencià; dels marcats «ampliat» se n'ha afegit l'estudi de la derivabilitat. Pots comprovar els resultats amb les eines dels apunts interactius.
\end{tcolorbox}

% =====================================================
\seccio{1}{Taxa de variació i derivada per definició}

\ex{1} Calcula la taxa de variació mitjana de les funcions següents en l'interval indicat:\begin{tasks}(3)
  \task $f(x)=x^3-x$ en $[0,2]$
  \task $g(x)=\ln x$ en $[1,e]$
  \task $h(x)=\sin x$ en $\left[0,\frac\pi2\right]$
\end{tasks}\recorda{$\text{TVM}[a,b]=\dfrac{f(b)-f(a)}{b-a}$: és el pendent de la recta secant. La derivada és el límit de la TVM quan l'interval es fa tan petit com vulguem.}

\ex{2} Llancem una pilota cap amunt i la seua altura, en metres, als $t$ segons és $s(t)=-5t^2+20t$. Calcula la velocitat mitjana en els dos primers segons i, aplicant la definició de derivada, la velocitat instantània en $t=1$.

\ex{2} Calcula, aplicant la definició de derivada:\begin{tasks}(2)
  \task $f'(2)$ si $f(x)=\dfrac{1}{x^2}$
  \task $f'(1)$ si $f(x)=\dfrac{x}{x+1}$
  \task $f'(0)$ si $f(x)=e^{3x}$
  \task $f'\!\left(\frac\pi2\right)$ si $f(x)=\sin x$
\end{tasks}\recorda{$f'(a)=\displaystyle\lim_{h\to0}\frac{f(a+h)-f(a)}{h}$. Amb exponencials i funcions trigonomètriques, usa els infinitèsims equivalents: $e^{u}-1\sim u$, \ $1-\cos u\sim\frac{u^2}{2}$, \ $\sin u\sim u$.}

\ex{3} Calcula la funció derivada aplicant la definició:\begin{tasks}(2)
  \task $f(x)=x^3$
  \task $f(x)=\dfrac{1}{\sqrt x}$
  \task $f(x)=\cos x$
  \task $f(x)=\ln(2x)$
\end{tasks}

\ex{2} \pau{Navarra, extraordinària 2023} Considera la funció $g(x)=-x^2+4x+3$. Calcula la derivada de $g(x)$ en el punt $x=1$ aplicant la definició de derivada.


% =====================================================
\seccio{2}{Rectes tangent i normal}

\ex{1} Escriu les equacions de les rectes tangent i normal a la gràfica de cada funció en el punt indicat:\begin{tasks}(2)
  \task $f(x)=x^3-3x+1$ en $x=2$
  \task $g(x)=\ln(x^2-3)$ en $x=2$
\end{tasks}\recorda{Tangent: $y-f(a)=f'(a)(x-a)$. \ Normal (perpendicular a la tangent): $y-f(a)=-\dfrac{1}{f'(a)}(x-a)$; si $f'(a)=0$, la normal és la recta vertical $x=a$.}

\ex{2} Troba l'equació de la recta tangent a la gràfica de cada funció en el punt indicat:\begin{tasks}(1)
  \task \pau{Múrcia, juliol 2020} $f(x)=\dfrac{1}{x+1}$ en $x=1$.
  \task \pau{Navarra, ordinària 2023} $f(x)=\sqrt{5+x^2}$ en $x=-2$.
  \task \pau{Navarra, ordinària 2022} $f(x)=\dfrac{x}{x^2-4}$ en $x=1$.
\end{tasks}

\ex{2} \pau{Balears, extraordinària 2021} Considera la funció $f(x)=\dfrac3x+8$.\begin{tasks}(1)
  \task Troba els punts de la gràfica en què la recta tangent és paral·lela a la recta $3x+4y+5=0$.
  \task Calcula les equacions de les rectes tangents en aquests punts.
\end{tasks}

\ex{2} \pau{Múrcia, setembre 2018} Donada la funció $f(x)=5x^3e^{2x}+\dfrac{1}{x^2+1}$, calcula $f'(0)$ i troba l'equació de la recta tangent a la gràfica en el punt $(0,1)$.

\ex{2} Donada $f(x)=x^3-6x^2+9x$, troba els punts on la tangent és horitzontal i l'equació de la recta tangent paral·lela a $y=-3x$.

\ex{3} Troba les rectes tangents a la paràbola $y=x^2$ que passen pel punt exterior $P(1,-3)$.\recorda{$P$ no és de la corba: escriu la tangent en un punt genèric $(a,a^2)$ i imposa que passe per $P$.}

\ex{2} Troba els paràmetres en cada cas:\begin{tasks}(1)
  \task \pau{Madrid, extraordinària 2022} $a$ i $b$ perquè $f(x)=ax+\dfrac bx$ verifique $f(2)=4$ i $f'(2)=0$.
  \task \pau{Múrcia, setembre 2017} $a$ i $b$ perquè $f(x)=\dfrac{ax+b}{x^2+1}$ complisca $f(0)=1$ i $f'(0)=1$.
  \task \pau{Andalusia, ordinària 2024} $a$ i $b$ perquè siga horitzontal la recta tangent a la gràfica de $h(x)=x^3+ax^2+3x+b$ en el punt $P(1,2)$.
\end{tasks}


% =====================================================
\seccio{3}{Regles de derivació i derivades successives}

\ex{1} Deriva i simplifica:\begin{tasks}(2)
  \task $f(x)=x^3\sqrt{x}-\dfrac{2}{\sqrt[3]{x}}$
  \task $f(x)=x\,e^x\sin x$
  \task $f(x)=\dfrac{x^2-1}{x^2+1}$
  \task $f(x)=\dfrac{x^3}{x-1}$
  \task $f(x)=\dfrac{1-\cos x}{\sin x}$
  \task $f(x)=e^x\left(x^2-2x+2\right)$
\end{tasks}\recorda{Escriu arrels i inverses com a potències. Producte de tres: $(fgh)'=f'gh+fg'h+fgh'$. Quocient: $\left(\frac fg\right)'=\frac{f'g-fg'}{g^2}$.}

\ex{1} \pau{Cantàbria, ordinària 2024} Considera la funció $f(x)=x\ln x$ amb $x>0$. Calcula la derivada de $f(x)$.

\ex{2} Comprova que $F(x)=x\arctan x-\frac12\ln(x^2+1)$ té per derivada $\arctan x$.

\ex{2} Derivades successives:\begin{tasks}(1)
  \task Calcula $f'''(x)$ si $f(x)=x^4-2x^3+x$.
  \task Troba la derivada $n$-èsima de $e^{-2x}$, de $x\,e^x$, de $\dfrac{1}{x+1}$ i de $\ln x$.
  \task Quina és la derivada 2027a de $\cos x$?
\end{tasks}\recorda{Calcula'n tres o quatre, busca el patró i escriu-lo en funció de $n$. Les derivades de $\sin x$ i $\cos x$ es repeteixen cada quatre.}

\ex{2} Troba $a$, $b$ i $c$ perquè $f(x)=ax^2+bx+c$ complisca $f(1)=2$, $f'(1)=3$ i $f''(1)=4$.

\ex{3} Comprova que la funció $y=e^{-x}\sin x$ verifica l'equació $y''+2y'+2y=0$.


% =====================================================
\seccio{4}{Regla de la cadena}

\ex{1} Deriva aquestes potències, arrels i exponencials:\begin{tasks}(3)
  \task $f(x)=(3x^2-1)^4$
  \task $f(x)=\sqrt{x^2-2x}$
  \task $f(x)=\dfrac{1}{(x^2+1)^3}$
  \task $f(x)=\sqrt[3]{2x+5}$
  \task $f(x)=e^{3x-1}$
  \task $f(x)=2^{x^2}$
\end{tasks}\recorda{$\left[f(g(x))\right]'=f'(g(x))\cdot g'(x)$: derivem la de fora, deixant dins $g(x)$, i multipliquem per la derivada de la de dins.}

\ex{2} Deriva aquestes funcions trigonomètriques i trigonomètriques inverses:\begin{tasks}(3)
  \task $f(x)=\sin^2(3x)$
  \task $f(x)=\cos(x^2)$
  \task $f(x)=\tan^3x$
  \task $f(x)=\arctan(2x)$
  \task $f(x)=\arcsin\sqrt x$
  \task $f(x)=\arccos e^{-x}$
\end{tasks}

\ex{2} Deriva i simplifica (aprofita les propietats dels logaritmes abans de derivar):\begin{tasks}(3)
  \task $f(x)=\ln(x^2-3x)$
  \task $f(x)=\ln\sqrt{\dfrac{1+x}{1-x}}$
  \task $f(x)=\log_2(x^3+1)$
  \task $f(x)=\ln(\cos x)$
  \task $f(x)=\ln\left(x+\sqrt{x^2+1}\right)$
  \task $f(x)=\ln\left(\tan\frac x2\right)$
\end{tasks}

\ex{2} Deriva i simplifica (hi ha més d'una regla en cada funció):\begin{tasks}(2)
  \task $f(x)=x\,e^{-x^2}$
  \task $f(x)=(x-1)^2(x+2)^3$
  \task $f(x)=\sqrt{x+\sqrt x}$
  \task $f(x)=e^{\sin x}\cos x$
  \task $f(x)=\sin(\ln x)+\cos(\ln x)$
\end{tasks}

\ex{2} \pau{Múrcia, setembre 2020} Calcula les derivades de les funcions següents:\begin{tasks}(2)
  \task $f(x)=(x^2-2)\ln x$
  \task $g(x)=e^{4x^3+2}$
\end{tasks}

\ex{2} \pau{Múrcia, juny 2018} Troba les derivades de les funcions següents:\begin{tasks}(2)
  \task $g(x)=x^2\ln(x^2)$
  \task $h(x)=e^{-3x+x^2}$
\end{tasks}

\ex{2} \pau{Andalusia, ordinària 2021} Calcula la derivada de les funcions següents:\begin{tasks}(2)
  \task $f(x)=\ln\left(\dfrac{x-1}{x+1}\right)$
  \task $g(x)=x^3\cdot e^{2x^2}$
\end{tasks}

\ex{2} \pau{Andalusia, ordinària (suplent) 2025} Calcula la derivada de les funcions següents:\begin{tasks}(2)
  \task $g(x)=\ln\left(\dfrac{x^2-1}{x^2+1}\right)$
  \task $h(x)=(2x-1)\,e^{x^2-x}$
\end{tasks}

\ex{3} \pau{Andalusia, ordinària 2024} Calcula la derivada de les funcions següents:\begin{tasks}(2)
  \task $f(x)=(x^2+2)^3\cdot e^{-2x}$
  \task $g(x)=\dfrac{\ln(1-x^3)}{(1-2x^2)^2}$
\end{tasks}

\ex{2} \pau{Múrcia, setembre 2018} Troba $a$ i $b$ perquè $f(x)=x^3\ln(2x+5)+ax+b$ complisca $f(0)=2$ i $f'(0)=1$.


% =====================================================
\seccio{5}{Derivació logarítmica}

\ex{2} Deriva aplicant la derivació logarítmica:\begin{tasks}(2)
  \task $y=x^x$
  \task $y=x^{\sin x}$
  \task $y=x^{\ln x}$
  \task $y=x^{1/x}$
  \task $y=(x^2+1)^x$
  \task $y=(\ln x)^x\ \ (x>1)$
  \task $y=(\cos x)^{\sin x}\ \ \left(0<x<\frac\pi2\right)$
  \task $y=x^{x^2}$
\end{tasks}\recorda{Prenem logaritmes: $\ln y=v\ln u$. Derivem els dos membres: $\dfrac{y'}{y}=v'\ln u+v\dfrac{u'}{u}$. Aïllem $y'$ i substituïm $y$.}

\ex{2} Deriva amb la derivació logarítmica (és més ràpid que amb les regles del producte i del quocient):\begin{tasks}(2)
  \task $y=\dfrac{x^3(x+1)^5}{\sqrt{x^2+4}}\ \ (x>0)$
  \task $y=\dfrac{(x-1)^2\sqrt[3]{x+2}}{(x+3)^4}\ \ (x>1)$
\end{tasks}

\ex{2} Considera la funció $f(x)=x^x$, amb $x>0$.\begin{tasks}(1)
  \task Troba les rectes tangents a la gràfica en $x=1$ i en $x=2$.
  \task En quin punt la tangent és horitzontal?
  \task Explica per què és incorrecte derivar-la com si fos una potència, $(x^x)'=x\cdot x^{x-1}$, o com si fos una exponencial, $(x^x)'=x^x\ln x$.
\end{tasks}


% =====================================================
\seccio{6}{Derivació implícita}

\ex{2} Calcula $y'$ en cada cas:\begin{tasks}(2)
  \task $x^2+y^2=4x$
  \task $x^3+y^3=3xy$
  \task $x^2y+xy^2=6$
  \task $e^{xy}=x+y$
  \task $\sqrt x+\sqrt y=3$
  \task $x=y+\ln y$
\end{tasks}\recorda{Derivem els dos membres respecte de $x$ pensant que $y=y(x)$: $(y^2)'=2y\,y'$, \ $(xy)'=y+x\,y'$, \ $(e^{y})'=e^{y}\,y'$. Després aïllem $y'$.}

\ex{2} Troba la recta tangent (i, en el primer cas, també la normal) a cada corba en el punt indicat. Comprova primer que el punt és de la corba.\begin{tasks}(1)
  \task $x^2+xy+y^2=7$ en $(2,1)$
  \task $x^3+y^3=9xy$ en $(2,4)$
  \task $\dfrac{x^2}{16}+\dfrac{y^2}{9}=1$ en $\left(\frac{16}{5},\frac95\right)$
  \task $y^2=x^3$ en $(1,1)$
\end{tasks}

\ex{3} En quins punts de la corba $x^2+xy+y^2=3$ la recta tangent és horitzontal? I vertical?

\ex{3} Calcula la derivada segona $y''$ i expressa-la en funció de $x$ i $y$ (o només de $y$):\begin{tasks}(3)
  \task $x^2+y^2=r^2$
  \task $x^2-y^2=1$
  \task $y^3+y=x$
\end{tasks}

\ex{2} Dedueix, derivant implícitament, les derivades de $y=\ln x$ (a partir de $e^y=x$), de $y=\arcsin x$ (a partir de $\sin y=x$) i de $y=\arctan x$ (a partir de $\tan y=x$).

\ex{3} Les dues tècniques alhora. El punt $(2,4)$ és de la corba $x^y=y^x$, ja que $2^4=4^2$. Troba el pendent de la tangent en aquest punt.\recorda{Pren primer logaritmes ($y\ln x=x\ln y$) i després deriva implícitament.}


% =====================================================
\seccio{7}{Continuïtat i derivabilitat}

\ex{2} Calcula les derivades laterals en el punt indicat i digues si la funció hi és derivable:\begin{tasks}(2)
  \task $f(x)=|x+1|$ en $x=-1$
  \task $f(x)=|x^2-9|$ en $x=3$
  \task $f(x)=x\,|x-1|$ en $x=1$
  \task $f(x)=\sqrt[3]{x^2}$ en $x=0$
  \task $f(x)=|\sin x|$ en $x=\pi$
  \task $f(x)=x^2|x|$ en $x=0$
\end{tasks}\recorda{$f'(a^\pm)=\displaystyle\lim_{h\to0^\pm}\frac{f(a+h)-f(a)}{h}$. És derivable si totes dues existeixen, són finites i són iguals.}

\ex{1} Raona si són vertaderes o falses. Si són falses, posa'n un contraexemple:\begin{tasks}(1)
  \task Si $f$ és contínua en $x=a$, aleshores és derivable en $x=a$.
  \task Si $f$ no és contínua en $x=a$, aleshores no és derivable en $x=a$.
  \task Si $f$ i $g$ no són derivables en $x=a$, $f+g$ tampoc no ho és.
  \task Si $f$ és derivable en $x=a$, $|f|$ també ho és.
\end{tasks}

\ex{2} \pau{Galícia, ordinària 2026} Dibuixa, quan siga possible, la gràfica d'una funció contínua però no derivable en $x=0$.

\ex{3} Siga $f(x)=x^2\sin\dfrac1x$ si $x\neq0$ i $f(0)=0$. Demostra, aplicant la definició, que $f$ és derivable en $x=0$. Passa el mateix amb $g(x)=x\sin\dfrac1x$, $g(0)=0$?\recorda{Si una funció tendeix a 0 i l'altra està fitada, el producte tendeix a 0.}


% =====================================================
\seccio{8}{Derivabilitat de funcions a trossos}

\ex{2} Estudia la continuïtat i la derivabilitat de les funcions següents:\begin{tasks}(2)
  \task $f(x)=\begin{cases}e^x&\text{si }x<0\\[3pt]x+1&\text{si }x\geq0\end{cases}$
  \task $g(x)=\begin{cases}3-x^2&\text{si }x<1\\[3pt]\dfrac2x&\text{si }x\geq1\end{cases}$
\end{tasks}\recorda{Primer la continuïtat. Després derivem tros per tros (sense l'igual en els punts de canvi) i, en cada punt de canvi on és contínua, comparem $f'(a^-)$ i $f'(a^+)$.}

\ex{2} \pau{Andalusia, extraordinària 2024} Estudia la continuïtat i la derivabilitat de\[f(x)=\begin{cases}-x^2+4x+3&\text{si }x<4\\[3pt]2x-5&\text{si }x\geq4\end{cases}\]

\ex{2} \pau{Andalusia, ordinària (suplent) 2025} Estudia la continuïtat i la derivabilitat de $f$ en el punt d'abscissa $x=-2$:\[f(x)=\begin{cases}10+\dfrac{5x}{2}&\text{si }x\leq-2\\[3pt]x^2+1&\text{si }-2<x<2\\[3pt]10-\dfrac{5x}{2}&\text{si }x\geq2\end{cases}\]

\ex{2} \pau{Andalusia, ordinària 2024} La velocitat del vent, en km/h, prevista per a cert dia en una zona de Sierra Nevada ve donada per\[v(t)=\begin{cases}t^2-8t+60&\text{si }0\leq t\leq10\\[3pt]-t^2+32t-140&\text{si }10<t\leq24\end{cases}\]on $t$ és el temps en hores. Comprova que la funció $v$ és contínua i derivable.

\ex{2} \pau{Navarra, ordinària 2024} Estudia la continuïtat (classificant les discontinuïtats) i la derivabilitat de\[f(x)=\begin{cases}-x^2-4x&\text{si }x<-1\\[3pt]4-x&\text{si }-1\leq x<1\\[3pt]x^2-6x+8&\text{si }x\geq1\end{cases}\]

\ex{2} \pau{Andalusia, extraordinària 2021} Estudia la continuïtat i la derivabilitat de\[f(x)=\begin{cases}2^{x+1}&\text{si }x<0\\[3pt]x^2-2x&\text{si }x\geq0\end{cases}\]

\ex{3} \pau{Madrid, extraordinària 2022 (ampliat)} Estudia la continuïtat i la derivabilitat en $x=0$ de\[f(x)=\begin{cases}x&\text{si }x\leq0\\[3pt]x\ln x&\text{si }x>0\end{cases}\]

\ex{2} \pau{Astúries, ordinària 2025 (ampliat)} Estudia la continuïtat i la derivabilitat de\[f(x)=\begin{cases}\dfrac x2&\text{si }0\leq x<4\\[3pt]3-(x-5)^2&\text{si }4\leq x\end{cases}\]

\ex{3} \pau{Andalusia, ordinària 2022 (ampliat)} Considera la funció contínua\[f(x)=\begin{cases}\dfrac1x&\text{si }x<-1\\[3pt]ax+b&\text{si }-1\leq x<1\\[3pt]\dfrac{x^2}{x+1}&\text{si }x\geq1\end{cases}\]\begin{tasks}(1)
  \task Calcula $a$ i $b$.
  \task Estudia la derivabilitat de $f$.
\end{tasks}

\ex{2} Estudia la derivabilitat de\begin{tasks}(2)
  \task $f(x)=\begin{cases}\sin x&\text{si }x\leq0\\[3pt]\ln(1+x)&\text{si }x>0\end{cases}$
  \task $g(x)=\begin{cases}\cos x&\text{si }x\leq0\\[3pt]1+x^2&\text{si }x>0\end{cases}$
\end{tasks}

\ex{3} \pau{Andalusia, extraordinària 2021} El nombre de diagnosticats per PCR a Andalusia, en milers de persones, s'aproxima per\[f(t)=\begin{cases}-t^2+2t-0{,}3&\text{si }0{,}2\leq t\leq1{,}8\\[3pt]0{,}1t-0{,}12&\text{si }1{,}8<t\leq5\\[3pt]-0{,}5t^2+8{,}3t-28{,}62&\text{si }5<t\leq10\end{cases}\]on $t$ és el temps en mesos. Estudia la continuïtat i la derivabilitat de $f$.

\ex{3} Escriu a trossos i estudia la derivabilitat de:\begin{tasks}(2)
  \task $f(x)=|x-1|+|x+1|$
  \task $g(x)=x\,|x-2|$
  \task $h(x)=e^{-|x|}$
  \task $k(x)=|x^2-5x+4|$
\end{tasks}


% =====================================================
\seccio{9}{Funcions a trossos amb paràmetres}

\ex{2} \pau{Andalusia, ordinària 2022} Determina $a$ i $b$ perquè $f$ siga contínua i derivable en tot el seu domini:\[f(x)=\begin{cases}6x-3&\text{si }x\leq1\\[3pt]ax^2+bx+2&\text{si }x>1\end{cases}\]\recorda{En cada punt de canvi: continuïtat ($\lim^-=\lim^+=f(a)$) i derivabilitat ($f'(a^-)=f'(a^+)$). Amb dos paràmetres, un sistema de dues equacions.}

\ex{2} \pau{Andalusia, extraordinària 2020} Calcula $a$ i $b$ perquè la funció siga contínua i derivable en el seu domini:\[f(x)=\begin{cases}2+\dfrac{a}{x-1}&\text{si }x<0\\[3pt]a+be^x&\text{si }x\geq0\end{cases}\]

\ex{2} Calcula els valors de $a$ perquè $f$ siga contínua i derivable:\begin{tasks}(1)
  \task \pau{Balears, ordinària 2022} \ $f(x)=\begin{cases}x^2-3x+2&\text{si }x<0\\[3pt]x^2+ax+2&\text{si }x\geq0\end{cases}$
  \task \pau{Balears, ordinària 2021} \ $f(x)=\begin{cases}x^2-3x+2&\text{si }x<0\\[3pt]e^{ax}+1&\text{si }x\geq0\end{cases}$
\end{tasks}

\ex{2} \pau{Balears, extraordinària 2021} Considera\[f(x)=\begin{cases}-4x+a&\text{si }x\leq-2\\[3pt]x^2-5&\text{si }-2<x<1\\[3pt]-7x+3&\text{si }x\geq1\end{cases}\]\begin{tasks}(1)
  \task Calcula els valors de $a$ perquè $f$ siga contínua.
  \task És $f$ derivable per a $a=1$?
\end{tasks}

\ex{3} \pau{Aragó, ordinària 2022} Sent $a$, $b$ paràmetres reals, es considera\[f(x)=\begin{cases}x^2+2\sqrt3&\text{si }x\leq0\\[3pt]\sqrt{ax+b}&\text{si }0<x\leq3\\[3pt]\dfrac72-\dfrac x6&\text{si }x>3\end{cases}\]\begin{tasks}(1)
  \task Determina el valor dels paràmetres perquè $f$ siga contínua.
  \task Per a aquests valors, analitza si $f$ és derivable en $x=0$ i en $x=3$.
\end{tasks}

\ex{2} \pau{Andalusia, ordinària (suplent) 2025} El nivell de concentració d'un alumne durant un examen de 5 hores ve donat per\[f(t)=\begin{cases}-t^2+2t+10&\text{si }0\leq t\leq2{,}5\\[3pt]t^2+at+b&\text{si }2{,}5<t\leq5\end{cases}\]on $t$ és el temps en hores. Amb quin nivell comença l'examen? Determina $a$ i $b$ perquè $f$ siga contínua i derivable en $t=2{,}5$.

\ex{3} \pau{Catalunya, ordinària 2026 (ampliat)} Determina els valors de $m$ que fan que $f$ siga contínua en tot el seu domini i, per a aquests valors, estudia'n la derivabilitat:\[f(x)=\begin{cases}5e^{2x}&\text{si }x\leq0\\[3pt](x+m)^2+1&\text{si }0<x<2\\[3pt]1&\text{si }2\leq x\end{cases}\]

\ex{2} \pau{Extremadura, extraordinària 2023 (ampliat)} Troba $a$ i $b$ perquè $f$ siga contínua en $x=1$ i la seua gràfica passe per $(-1,5)$. És derivable en $x=1$?\[f(x)=\begin{cases}2x^2+ax+b&\text{si }x\leq1\\[3pt]\ln x&\text{si }x>1\end{cases}\]

\ex{2} \pau{Aragó, extraordinària 2021 (ampliat)} Calcula els valors de $a\neq0$ perquè $f$ siga contínua i estudia si, per a algun d'aquests valors, és derivable en $x=1$:\[f(x)=\begin{cases}5-ax^2&\text{si }x\leq1\\[3pt]\dfrac{6}{ax}&\text{si }x>1\end{cases}\]

\ex{2} Troba els paràmetres perquè cada funció siga derivable en el seu domini:\begin{tasks}(2)
  \task $f(x)=\begin{cases}ax^2+bx&\text{si }x\leq1\\[3pt]\ln x&\text{si }x>1\end{cases}$
  \task $g(x)=\begin{cases}ax^2+b&\text{si }x<1\\[3pt]\dfrac1x&\text{si }x\geq1\end{cases}$
  \task $h(x)=\begin{cases}e^{ax}+b&\text{si }x\leq0\\[3pt]\sin2x+3&\text{si }x>0\end{cases}$
  \task $k(x)=\begin{cases}a+bx&\text{si }x\leq0\\[3pt]\arctan x&\text{si }x>0\end{cases}$
\end{tasks}

\vfill
\begin{center}\color{muted}\small\faCheckCircle[regular]\ \ Ànim! Comprova les teues derivades amb la calculadora dels apunts interactius.\end{center}

\end{document}
